% LSPV-5 REPORT::PRIMEPH1::2026-10-02T09:45:00-07:00
\documentclass[11pt]{article}
\usepackage[margin=.88in]{geometry}
\usepackage{amsmath,amssymb}
\begin{document}
\section*{PRIMEPH1: polynomial-hierarchy closure of prime-bound provenance search}
\textbf{STATUS: PROVED/EXACT closure theorem + COMPUTATIONAL finite control. P versus NP remains OPEN.}

\section*{CORPUS REVIEWED}
The current repository tree/history was surveyed. This run reviewed root \texttt{main.tex};
AGENTS/CODEX/pnp-research governance; TASK-001--007; verified MAIN 0858 and relevant predecessors;
PRIME-LQP-1 and PRIMECRT1; MATRIXCERT1; EXCESS1, EXCESSDUAL1, and concurrent TURANEX1;
RAMCOMP1; predecessor TRACECLOSE1; current GROK/CLAUDE and independent-review artifacts;
HEC/VSC-HEC/G-PPC-HEC/CRS-PPC/CCSPC/RCCN/Search--G; and active experiment/result files.
Predecessor LSPV-4 TRACECLOSE1 verifies, so no recovery write is needed.

\section*{LITERATURE PREFLIGHT}
Stockmeyer and Wrathall give the polynomial hierarchy and polynomially bounded quantifier
characterization. Megiddo--Papadimitriou introduced TFNP. Kleinberg--Korten--Mitropolsky--
Papadimitriou and Fleming--Imrek--Marciot study total search higher in PH, including
TF\(\Sigma_2\) and unique-solution variants. Function-PH work also studies function analogues
of PH levels. These are relevant taxonomies, but under the counterfactual \(P=NP\), PH collapses
to P. Provenance-semiring/weighted-rewriting results validate derivation annotations but do not
force arbitrary static SAT deciders to preserve a chosen annotation.

\section*{PROVED/EXACT: PRIME-PH-CLOSE1}
Let \(R(x,w)\in PH\) be polynomially balanced, so valid witnesses satisfy
\[
|w|\le p(|x|).
\]
Assume \(P=NP\). Then on
\[
D_R=\{x:\exists w\ R(x,w)\}
\]
there is an FP selector \(S\) with \(R(x,S(x))\).

For a prefix \(u\), define
\[
E_R(x,u)\iff \exists v\ R(x,uv),\qquad |uv|\le p(|x|).
\]
Since \(R\in PH\), adjoining one bounded existential block keeps \(E_R\in PH\).
Under \(P=NP\), \(PH=P\), so \(E_R\) is polynomial-time decidable. Query successively
whether \(E_R(x,u0)\) holds; append 0 if yes and 1 otherwise. At most \(p(|x|)\) queries
recover a valid witness. Thus \(S\in FP\).

If \(g(x,w)\) is any polynomial-time provenance annotation of polynomial output length, then
\[
H(x)=g(x,S(x))\in FP.
\]
This includes positional primes, CRT residues, and source/witness/gate event labels.

\paragraph{Scope.}
This is qualitative polynomial closure, not a tight reconstruction exponent. RAMCOMP1 remains the
appropriate warning that qualitative \(PH=P\) is not a free quantitative compiler.

\section*{Corollary: uniqueness and canonicalization}
A unique valid witness is still selected by PRIME-PH-CLOSE1. With multiple witnesses the prefix
rule selects the lexicographically first valid witness. Finite-PH uniqueness/minimality conditions
remain within PH after a fixed number of bounded quantifier blocks. Therefore unique or canonical
prime-bound provenance does not by itself evade reacquisition under \(P=NP\).

\section*{Corollary: current reconstruction prefix}
For the current dependency
\[
\exists\sigma\,\forall u\,\exists z\,\forall y\,
B(x,\sigma,u,z,y),\qquad B\in P,
\]
with polynomially bounded strings, set
\[
Good(x,\sigma)\iff
\forall u\,\exists z\,\forall y\,B(x,\sigma,u,z,y).
\]
Then \(Good\in\Pi_3^P\), while the prefix-extension predicate for \(\sigma\) is in
\(\Sigma_4^P\). Under \(P=NP\), both collapse to P, so a valid \(\sigma\), when one
exists, is selectable in FP. Any polynomial-time prime annotation of that selected
\(\sigma\) is FP.

This does not solve reconstruction. Exact source typing, model matching, and quantitative compiler
charges remain OPEN. It does rule out fixed finite alternation as a generic provenance-based source
of superpolynomial reacquisition.

\section*{COMPUTATIONAL control: PRIMEPH1}
Artifacts:
\[
\texttt{experiments/prime\_lqp/prime\_ph\_closure\_test.py},
\qquad
\texttt{results/raw/prime\_lqp/2026-10-02\_LSPV5\_PRIMEPH1.json}.
\]
The exhaustive finite control uses
\[
\exists\sigma_5\forall u_3\exists z_3\forall y_2 R.
\]
For a unique instance it recovers \(\sigma=10110\) in five prefix queries and prime annotation
161700. For a nonunique instance with valid witnesses
\[
\{00011,10110,11111\},
\]
it recovers lexicographically first \(00011\) and annotation 177870.
Status: \textbf{COMPUTATIONAL}. This tests the implementation pattern, not \(P=NP\).

\section*{Hostile classification}
A prime/LQP consumer now falls into one of four buckets:
\begin{enumerate}
\item semantic quotient: provenance is irrelevant;
\item deterministic polynomial transcript functional: closed by TRACECLOSE1;
\item polynomially balanced PH replacement relation: closed under \(P=NP\) by PRIME-PH-CLOSE1;
\item outside-PH or exact-original residue: not closed by the generic arguments, but still requiring
an arbitrary-algorithm necessity theorem and survival of reread, RECOMP, semantic/algebraic oracle,
wrapper, and representation attacks.
\end{enumerate}
Moving to PSPACE- or exponential-verification provenance is not automatically progress; absent a
model-preserving necessity theorem it merely imports the stronger computation into the definition.

\section*{DUPLICATION / CONTRADICTION AUDIT}
PRIME-FNP-CLOSE1 is the special case with a P-time verifier. PRIME-PH-CLOSE1 broadens the project
closure to any fixed PH-verifiable replacement relation. It is compatible with RAMCOMP1 because it
claims polynomial closure only, not a favorable exponent. It is compatible with ISOCAN1 and
TRACECLOSE1. Concurrent TURANEX1 is independent and strengthens only the structural support/
generalized-weight diagnostics.

\section*{WHAT CHANGED RELATIVE TO ROOT MAIN}
The provenance barrier should now read:
\[
\boxed{\text{semantic factoring / deterministic replay / polynomially balanced PH replacement search}.}
\]
The prior surviving condition stating only not-NP-verifiable replacement is stale and too weak.
The current four-alternation reconstruction relation is explicitly covered qualitatively.

\section*{ADOPTED / REJECTED / RETIRED}
\textbf{ADOPT:} PRIMECRT1; event primes as instrumentation; TRACECLOSE1/REACQ1;
PH/TFPH literature as taxonomy; TURANEX1 in structural scope.

\textbf{REJECT AS TRANSFER:} TFPH, provenance-semiring, weighted-rewriting, or function-PH results
do not themselves force unrestricted SAT runtime.

\textbf{FALSIFIED/RETIRED AS GENERIC HINGES:} finite-alternation prime provenance; unique or
canonical PH-verifiable prime-bound witnesses; prime magnitude/factorization; deterministic
polynomial trace tags.

\section*{STALE OR OVERSTATED CLAIMS FOUND}
Any claim that fixed finite alternation alone takes a polynomial-length SAME\_REALIZATION
provenance relation outside reacquisition under \(P=NP\) is overstated. The correct exclusion
must be outside polynomially balanced PH replacement search, not merely outside FNP.

\section*{Barrier audit}
The closure theorem relativizes and supplies no nonalgebrizing hinge. DEAD/COPY/DELAY alter
provenance without semantics; INTERP/RECOMP handle deterministic histories; padding may inflate
codes. Horn, dual-Horn, 2-SAT, affine/Tseitin, bounded-treewidth, renaming, and independent-product
controls remain compatible. No circuit property is asserted.

\section*{Single next provenance target}
Continue PRIME-LQP only if a source-matched later restriction \(D\) on two SU/vMS/RMV histories
with the same semantic state satisfies
\[
D\not\text{ factor through semantics},\qquad
D\notin\text{ polynomially balanced PH replacement search},
\]
and is not deterministic-trace replayable, while a model-preserving theorem forces every relevant
static SAT decider to realize \(D\). Otherwise prime provenance should remain instrumentation.

\begin{thebibliography}{9}
\bibitem{Stockmeyer76} L. J. Stockmeyer, The Polynomial-Time Hierarchy, Theoretical Computer Science 3(1), 1976, 1--22.
\bibitem{Wrathall76} C. Wrathall, Complete Sets and the Polynomial-Time Hierarchy, Theoretical Computer Science 3(1), 1976, 23--33.
\bibitem{MP91} N. Megiddo, C. H. Papadimitriou, On total functions, existence theorems and computational complexity, Theoretical Computer Science 81(2), 1991, 317--324.
\bibitem{KKMP21} R. Kleinberg, O. Korten, D. Mitropolsky, C. Papadimitriou, Total Functions in the Polynomial Hierarchy, ITCS 2021, LIPIcs 185:44.
\bibitem{FIM25} N. Fleming, D. Imrek, C. Marciot, Provably Total Functions in the Polynomial Hierarchy, CCC 2025, LIPIcs 339:28.
\bibitem{FunPH22} The polynomial hierarchy of functions and its levels, Theoretical Computer Science 900 (2022), 25--34.
\end{thebibliography}
\end{document}
% END_OF_VERIFIED_REPORT::LSPV-5::PRIMEPH1
